\begin{pastexam}[label={px:ntu-110-comm-dsa-37}]{臺大電信所 110 學年度（B）第 37 題}
（複選；每題 2.5 分，每答錯一個選項倒扣 0.5 分至該題 0 分止）

Which of the following statements are true?
\begin{enumerate}[label=\Alph*.]
\item Stack implements a FIFO (first-in-first-out) policy.
\item Queue implements a LIFO (last-in-first-out) policy.
\item Stack can be realized by deque (double-ended queue).
\item The Find() operation in a doubly linked list takes O(n) time.
\item The Enqueue() operation in a queue takes O(1) time.
\end{enumerate}
\end{pastexam}

\begin{pastsoln}
答案是 C、D、E。
\begin{itemize}
\item A 錯：堆疊是後進先出（LIFO），先進先出的是佇列（資結書第~6.1~節）。
\item B 錯：佇列是先進先出（FIFO），反過來了（資結書第~6.2~節）。
\item C 對：雙端佇列兩端都能進出，只用其中一端（例如只用 push-back、pop-back）就是堆疊；只從一端進、
另一端出就是佇列（資結書第~6.2.6~節）。
\item D 對：鏈結串列不能隨機存取，找一個值要從頭一個一個比，最壞 $O(n)$；雙向串列多了往回的指標，但
搜尋還是要走過每個節點（資結書第~5.2~節）。
\item E 對：用環狀陣列（容量足夠時）或帶尾指標的鏈結串列實作，加入只動尾端，$O(1)$（資結書第~6.2.2~節）。
\end{itemize}
\end{pastsoln}
